\section{Experimental Setup}

\subsection{Data}
Table~\ref{tab:data} lists the tasks. \fpb{}, \fiqa{} and Headlines use the FLARE
splits~\cite{xie2023pixiu}. \tfns{} has no public test set; its official validation split is
our test set and a stratified 10\% of the training data is held out for validation. Headlines
asks several yes/no questions about the same headline (price up/down, past/future news, \ldots);
we remove from training every headline that appears in validation or test, so that no headline
is seen under a different question. The two large sources are subsampled with a fixed seed
(Table~\ref{tab:data}); all splits are stored as files in the repository.
SST-2~\cite{socher2013sst} (872 validation sentences) is used only as the general-domain
reference for the routing analysis and is never trained on.

\begin{table}[t]
\centering
\caption{Tasks and split sizes after subsampling.}
\label{tab:data}
\begin{tabular}{l l c r r r}
\toprule
Task & Source & Cls. & Train & Val & Test \\
\midrule
\fpb{} & news sentences~\cite{malo2014fpb} & 3 & 3100 & 776 & 970 \\
\fiqa{} & posts, microblogs~\cite{maia2018fiqa} & 3 & 750 & 188 & 235 \\
\tfns{} & tweets~\cite{tfns2022} & 3 & 2500 & 500 & 2388 \\
Headlines & gold news~\cite{sinha2021headlines} & 2 & 2500 & 500 & 2000 \\
SST-2 & movie reviews~\cite{socher2013sst} & 2 & -- & -- & 872 \\
\bottomrule
\end{tabular}
\end{table}

\subsection{Models}
The main model is Granite-3.1-1B-A400M-base~\cite{granite2024}: 24 layers, hidden size 1024,
32 experts per layer with intermediate size 512, top-8 routing, 1.33B parameters in total and
about 0.4B active per token. As a dense reference we use Qwen2.5-0.5B~\cite{qwen2025qwen25}
(0.49B parameters, all active) with the same attention LoRA.

\subsection{Training}
One adapter is trained on the union of the four training sets for one epoch (553 steps of
batch size 16) with AdamW~\cite{loshchilov2019adamw}, no weight decay, 6\% linear warm-up and
cosine decay, gradient clipping at 1.0 and maximum prompt length 192 tokens (longer prompts are
truncated from the left so that the answer cue is kept). The learning rate
is $2\cdot10^{-4}$ for LoRA, $5\cdot10^{-4}$ for the router and $3\cdot10^{-5}$ for the
full-rank ESFT deltas; LoRA uses $\alpha=2r$. The frozen backbone runs in fp16 with dynamic
loss scaling and all trainable tensors are kept in fp32. The validation set is evaluated every 200 steps and at
the end, and the adapter with the best mean validation macro-F1 is used for testing. No
hyper-parameter search was performed: the values follow common practice for each method and
are identical across tasks. Every configuration is trained with three seeds (0, 1, 2), which
change the data order, the LoRA initialisation and, for ESFT-Random, the chosen experts. All
runs use a single NVIDIA T4 (16\,GB) on Kaggle.

\textbf{Padding in fp16.} Batches are right-padded. With Granite's default padding token
(\texttt{<|endoftext|>}) the hidden states of padded positions overflow the fp16 range in the
last layers. The prediction is read at the last real token, which never attends to padding, so
inference is unaffected, but in the backward pass the zero gradient at these positions meets
an infinite activation and produces NaN gradients in almost every step. Padding with an ordinary token removes
the problem; a unit test checks that a padded example gets exactly the same logits as the
same example alone.

\subsection{Metrics}
We report macro-F1 (primary, because label distributions are skewed in opposite directions
across tasks) and accuracy per task, and their unweighted mean over the four financial tasks.
Efficiency is measured by the number of trainable parameters, training time, peak accelerator
memory and inference latency (batch 16, T4). Frozen-model baselines are deterministic and are
evaluated once.
